\section{Local fourth moments at a prescribed physical endpoint}
\label{sec:v36-local-bridge-moments}
A positive microscopic denominator is not by itself a tightness
estimate for a conditioned path. We prove a local fourth-moment bound
before conditioning. Its factor $m^{-3/2}$ exactly pays that denominator.
The proof differentiates only a fixed block in the collision product;
it never pulls a path indicator through an induced iterate.

Use a compact family satisfying \textup{(H1)--(H7)}. Write
$f=f^{\rm c}_\lambda=(\kappa,\tau-\bar\tau_\lambda)$,
$Q_z=Q_{\lambda,z}$, and
\[
 V_{r,l}=\sum_{j=r}^{r+l-1}f\circ T_\lambda^j,
 \qquad 0\le r\le r+l\le m.
\]
The time sum $S_m$ below is the uncentered sum of roofs, whereas
$V_{r,l}$ is centered. Norms and constants are taken on a local common
space and then made uniform by a finite parameter cover.

\begin{lemma}[A differentiated collision block]
\label{lem:v36-differentiated-block}
There are $\epsilon,c,C>0$ such that for every unit real vector $v$,
$l\ge1$, and real $|\omega|\le\epsilon$,
\begin{align}
 \|Q_\omega^j\|&\le Ce^{-cj|\omega|^2}\quad(j\ge0),
                                            \label{eq:v36-power-Gaussian}\\
 \left\|\left.\partial_t^4Q_{\omega+tv}^{\,l}\right|_{t=0}\right\|
 &\le C(l^2+l^4|\omega|^4)e^{-cl|\omega|^2}.
                                            \label{eq:v36-fourth-block}
\end{align}
For each fixed $B<\infty$, on
$\omega\in\T^2\times[-B,B]$, $|\omega|\ge\epsilon$, there are
$C_B<\infty$ and $\rho_B<1$ such that
\begin{equation}\label{eq:v36-away-block}
 \left\|Q_\omega^{m-r-l}
   \left.\partial_t^4Q_{\omega+tv}^{\,l}\right|_{t=0}
                   Q_\omega^r\right\|
       \le C_B(1+l)^4\rho_B^{m}.
\end{equation}
The constants in the last assertion are chosen for fixed $B$.
\end{lemma}
\begin{proof}
On a fixed complex neighborhood of zero the spectral decomposition is
$Q_z^l=e^{l\psi(z)}\Pi(z)+E_l(z)$, where the projection and logarithm
are analytic, $\|E_l(z)\|\le C\rho^l$, and
\[
 \psi(0)=0,\quad \psi'(0)=0,\quad
 \psi''(0)=-\Sigma_\lambda,
 \qquad \operatorname{Re}\psi(\omega)\le-c|\omega|^2
\]
for real small $\omega$. Uniform ellipticity and shrinking the
neighborhood give the last inequality. The complementary powers can
be bounded on a slightly larger complex neighborhood, so Cauchy's
inequality also gives $\|\partial_t^jE_l(\omega+tv)|_{t=0}\|
\le C_j\rho_1^l$ for $j\le4$. The same decomposition proves
\eqref{eq:v36-power-Gaussian}, after decreasing $c$ to include the
complementary power.

All fixed derivatives of $\Pi$ and $\psi$ are bounded, and
$|\partial_v\psi(\omega)|\le C|\omega|$. The fourth derivative of
$e^{l\psi}$ is $e^{l\psi}$ times
\[
 l\psi^{(4)}+4l^2\psi'\psi^{(3)}
 +3l^2(\psi'')^2+6l^3(\psi')^2\psi''+l^4(\psi')^4.
\]
Leibniz's rule includes the lower derivatives multiplied by derivatives
of $\Pi$. Their absolute sum is bounded by
$C(l^2+l^4|\omega|^4)$, since
$l^3|\omega|^2\le(l^2+l^4|\omega|^4)/2$ and $l\ge1$.
The complementary derivative is absorbed by the same Gaussian bound.
This proves \eqref{eq:v36-fourth-block}.

Away from zero, every unmarked real power has bound $C_B\rho_B^j$.
Differentiate the product of the $l$ one-step operators four times.
Each resulting word contains at most four differentiated one-step
operators, uniformly bounded on the fixed band, and at most five
undifferentiated blocks. Their lengths together with the two exterior
blocks sum to at least $m-4$. There are at most $C(1+l)^4$ such terms.
Multiplying their spectral bounds and absorbing the fixed missing
powers proves \eqref{eq:v36-away-block}. This argument does not infer
a derivative bound from a complex growing-band estimate.
\end{proof}

\begin{proposition}[A local fourth-moment bound]
\label{prop:v36-local-fourth}
For each fixed bounded interval $J\subset\R$ there is $C_J$ such that,
for $m\ge1$, $1\le l\le m$, every admissible $r$, and all $k\in\Z^2$,
$t\in\R$,
\begin{equation}\label{eq:v36-local-fourth}
 \int |V_{r,l}|^4\one_{\{K_m=k,\ S_m-t\in J\}}\,d\nu
                       \le C_J l^2m^{-3/2}.
\end{equation}
The bound is uniform in $\lambda$ and in the location of the block.
No central restriction on $(k,t)$ is needed in this upper bound.
\end{proposition}
\begin{proof}
Choose once a nonnegative integrable majorant $q_J\ge\one_J$ with
integrable Fourier transform supported in a fixed interval $[-B,B]$,
as in the positive-envelope construction of
Theorem~\ref{thm:two-endpoint-collision-LLT}. It is fixed independently
of $m,l,k,t,\lambda$. For a unit coordinate vector $v$, positivity
bounds the corresponding scalar fourth moment by the same integral
with $q_J(S_m-t)$ in place of $\one_J$.

Torus orthogonality and Fourier inversion represent this integral as
\begin{align*}
 \frac1{(2\pi)^3}\int e^{-i[u\cdot k+b(t-m\bar\tau_\lambda)]}
 \widehat q_J(-b)\,
 \ell\left(Q_\omega^{m-r-l}
       \left.\partial_z^4Q_{\omega+zv}^{\,l}\right|_{z=0}
                         Q_\omega^r\nu\right)\,du\,db.
\end{align*}
There is no minus sign from four differentiations, since $i^4=1$.
The identity follows first for finite collision sums, which are
bounded, and then by ordinary dominated Fourier integration on this
fixed band. The exterior powers occur on the correct sides of the
marked block.

The previous lemma bounds the integrand near zero by
$C(l^2+l^4|\omega|^4)e^{-cm|\omega|^2}$.
Its three-dimensional integral is at most
\[
 C\bigl(l^2m^{-3/2}+l^4m^{-7/2}\bigr)
                      \le 2C l^2m^{-3/2},
\]
since $l\le m$. The remaining integral is bounded by
$C_B(1+l)^4\rho_B^m$, which is at most $C l^2m^{-3/2}$
for $1\le l\le m$. The Fourier phase has modulus one, explaining
uniformity in the target. Finally
$|V|^4\le3\sum_{i=1}^3|V_i|^4$ proves the vector estimate.
\end{proof}

Let $A_{t,\lambda}(k,m)=\{W_\lambda(t)=k,C_\lambda(t)=m\}$, and let
$Z_m$ be the polygonal collision path with
\[
 Z_m(j/m)=m^{-1/2}\sum_{i<j}f_\lambda^{\rm c}\circ T_\lambda^i,
                           \qquad 0\le j\le m.
\]
A stationary initial state is $(x,a)$, with $x$ the last outgoing
collision and $a$ its actual age. Thus $Z_m$ is a function of that
same source trajectory, not of a new process.

\begin{corollary}[Conditional tightness on the microscopic event]
\label{cor:v36-conditional-tightness}
For every fixed central compact set $|\xi_{t,\lambda}(k,m)|\le M$,
and all sufficiently large $t$,
\begin{equation}\label{eq:v36-conditional-fourth}
 \mathbb E_{\mu_\lambda}\left[
  |Z_m((r+l)/m)-Z_m(r/m)|^4\mid A_{t,\lambda}(k,m)\right]
                           \le C_M(l/m)^2.
\end{equation}
The conditional polygonal path laws are uniformly tight in
$C([0,1],\R^3)$. The same conclusion holds with nonnegative bounded
regular endpoint selectors whose product of stationary means has a
uniform positive lower bound.
\end{corollary}
\begin{proof}
The exact age integral for the numerator of a nonnegative function of
$x$ is the finite sum over initial and terminal cell offsets $(d,e)$.
Each term has label $K_m=k+d-e$ and overlap
$F_\lambda^{d,e}(x,S_m-t,T^mx)$. This overlap is at most
$\tau_+\one_{[-\tau_+,\tau_+]}(S_m-t)$. Apply
Proposition~\ref{prop:v36-local-fourth} to each term. The source
factor $1/\bar\tau_\lambda$ is uniformly bounded. The resulting
unnormalized fourth moment of $V_{r,l}$ is at most
$Cl^2m^{-3/2}$. The actual event probability is at least
$c_Mt^{-3/2}$, and $m\asymp t$ on the specified central set.
Division and the factor $m^{-2}$ give
\eqref{eq:v36-conditional-fourth}.

At arbitrary times, split a polygonal increment into its two partial
grid intervals and the complete middle interval. On a single grid
interval a fraction $\theta$ multiplies its increment by $\theta$;
its fourth moment is at most $C\theta^4m^{-2}$, hence at most the
square of the elapsed time times $C$. The elementary inequality
$|u+v+w|^4\le27(|u|^4+|v|^4+|w|^4)$ gives the same
$C|s-t|^2$ bound across grid intervals. The continuity tightness
criterion, or its dyadic proof, now gives uniform tightness, since
$Z_m(0)=0$. For selected probabilities the numerator is bounded by
the selector supremum product times the one just estimated. The
selected local law supplies its positive denominator. No mixing
estimate for a long path indicator is used.
\end{proof}
